% ============================================================================
%  TP Final - Simulación de Sistemas (72.25) - ITBA - 2026 Q1 - Grupo G01S2
%  Presentación (Beamer). Compilar: pdflatex dos veces.
%  Figuras en ../figures/ ; fotogramas hero en ../data_anim/.
%  Los fotogramas van como IMAGEN FIJA (política del README: en el PDF entregable no van
%  animaciones ni enlaces run: a archivos locales). Cuando estén los enlaces web de las
%  animaciones (YouTube/Vimeo), reemplazar \videohero por \videolink (ver preámbulo).
%  Duración objetivo: ~12-15 min. Diapositivas numeradas.
% ============================================================================
\documentclass[10pt,aspectratio=169]{beamer}

\usepackage[utf8]{inputenc}
\usepackage[T1]{fontenc}
% babel-spanish no está en este TeX Live (Debian); se omite para compilar de forma portable. Los
% acentos salen por inputenc/fontenc y las diapositivas no usan rótulos automáticos en inglés.
\usepackage{lmodern}
\usepackage{amsmath,amssymb}
\usepackage{graphicx}
\usepackage{booktabs}
\renewcommand{\figurename}{Figura}
\renewcommand{\tablename}{Tabla}

\usetheme{Boadilla}
\usecolortheme{whale}
\setbeamertemplate{navigation symbols}{}
\setbeamertemplate{caption}{\raggedright\insertcaption\par}
\setbeamerfont{footnote}{size=\tiny}
\graphicspath{{../figures/}{../data_anim/}}

% Separadores de sección: una diapositiva-título breve al abrir cada sección.
\setbeamertemplate{section in toc}[sections numbered]
\AtBeginSection[]{%
  \begin{frame}[plain,noframenumbering]\centering\vfill
    {\usebeamercolor[fg]{structure}\structure{\Large\textbf{\insertsectionhead}}}\vfill
  \end{frame}}

\newcommand{\unit}[1]{\,\mathrm{#1}}
% IMAGEN FIJA de un fotograma representativo (sin enlace run: a archivos locales). Args: base, ancho.
\newcommand{\videohero}[2]{\includegraphics[width=#2]{#1_fotograma.png}}
% Cuando existan los enlaces web, usar en su lugar: \videolink{base}{ancho}{URL}
\newcommand{\videolink}[3]{\href{#3}{\includegraphics[width=#2]{#1_fotograma.png}}}
\newcommand{\playnota}{{\scriptsize Fotograma representativo; la animación completa se publicará con enlace web.}}

\title[Diagrama fundamental de VDV]{Diagrama fundamental de vehículos dirigidos por vibración (VDV)}
\subtitle{Reproducción de tendencias con un autómata celular de Nagel--Schreckenberg con contacto}
\author[Grupo G01S2]{Nicolás V. Arias (Leg.~62272) \and Keoni L. Dubovitsky (Leg.~62815) \and Ian F. Tognetti (Leg.~61215)}
\institute[ITBA]{72.25 Simulación de Sistemas --- ITBA}
\date{2026, primer cuatrimestre}

\begin{document}

\begin{frame}[plain]\titlepage\end{frame}

% ============================================================================
\section{Problema y modelo}
% ============================================================================

% ---- El sistema real --------------------------------------------------------
\begin{frame}{El sistema real: VDV}
\begin{itemize}
  \item Robots Hexbug Nano ($44\times15\times18\unit{mm}$): la vibración rectificada por cerdas
  asimétricas los impulsa.
  \item Canal \textbf{circular} de $\approx1313\unit{mm}$: sistema \textbf{1D}, hasta $30$ vehículos.
  \item Interactúan \textbf{solo por contacto} (sin percepción remota: no ``ven'' al de adelante).
  \item Velocidades libres heterogéneas, $\sim\mathcal{U}[90,120]\unit{mm/s}$.
\end{itemize}
\vfill
{\footnotesize Aísla la parte del diagrama fundamental debida \textbf{solo a colisiones físicas}.}
\end{frame}

% ---- Lo que mide el experimento ---------------------------------------------
\begin{frame}{Lo que mide el experimento (Patterson \& Parisi, 2023)}
\begin{itemize}
  \item Velocidad media \textbf{casi constante} a densidad baja/media.
  \item \textbf{Cae 25--40\%} cerca de la densidad de contacto ($1/(44\unit{mm})$).
  \item El \textbf{orden de inserción} importa: el aleatorio queda \emph{acotado} entre creciente y decreciente.
  \item A saturación ($N=30$) \textbf{convergen}, por debajo del más lento ($90\unit{mm/s}$).
\end{itemize}
\vfill
{\footnotesize \textbf{Objetivo:} reproducir estas \emph{tendencias} con un autómata celular de bajo costo.}
\end{frame}

% ---- Modelo -----------------------------------------------------------------
\begin{frame}{Modelo: Nagel--Schreckenberg con contacto}
\begin{columns}[T]
\begin{column}{0.54\textwidth}
Ruta \textbf{periódica} 1D de $L$ celdas; vehículo de $\ell$ celdas; $v_i\in\{0,\dots,v_{\max,i}\}$.
Paso sincrónico, orden $R1\to R3\to R2\to R4$:
\begin{itemize}\small
  \item[R1] acelerar: $v\leftarrow\min(v+1,v_{\max})$
  \item[R3] frenar (prob. $p$): $v\leftarrow\max(0,v-1)$
  \item[R2] proyectar contactos (sin solapar)
  \item[R4] mover: $x\leftarrow(x+d)\bmod L$
\end{itemize}
\end{column}
\begin{column}{0.42\textwidth}
\footnotesize
Frenar \emph{antes} de proyectar $\Rightarrow$ nunca hay solapamiento.
\medskip

\textbf{Contacto puro:} avanza libre hasta el líder; a contacto hereda su velocidad.
\[ d_i=\min(\hat v_i,\; g_i+d_{i+1}) \]
\end{column}
\end{columns}
\end{frame}

% ---- Dos variantes ----------------------------------------------------------
\begin{frame}{Dos variantes de la Regla 2}
\begin{columns}[T]
\begin{column}{0.5\textwidth}
\textbf{Contacto puro (A) --- oficial}
\begin{itemize}\small
  \item Sin anticipación: flujo libre hasta el contacto.
  \item Al alcanzar al líder, hereda su velocidad.
  \item Es la física del experimento.
\end{itemize}
\end{column}
\begin{column}{0.5\textwidth}
\textbf{Clásica (B) --- validación}
\begin{itemize}\small
  \item NaSch de libro: $v\leftarrow\min(v,g)$.
  \item Solo para \textbf{validar} contra la solución analítica.
  \item Homogénea, $p=0$, $\ell=1$.
\end{itemize}
\end{column}
\end{columns}
\end{frame}

% ---- Calibración ------------------------------------------------------------
\begin{frame}{Calibración a la geometría del experimento}
\centering
\begin{tabular}{lll}
\toprule
Cantidad & Símbolo & Valor \\
\midrule
Paso espacial & $\Delta x$ & $0{,}25\unit{mm}$ \\
Paso temporal & $\Delta t$ & $1/24\unit{s}$ (24 cuadros/s) \\
Cuanto de velocidad & $\Delta v=\Delta x/\Delta t$ & $6\unit{mm/s}$ \\
Largo de vehículo & $\ell$ & $176$ celdas $=44\unit{mm}$ \\
Largo de pista & $L$ & $5280$ celdas $=1320\unit{mm}=30\,\ell$ \\
\bottomrule
\end{tabular}
\vfill
{\footnotesize $L=30\,\ell$ hace que $N=30$ cierre \emph{exacto} a contacto ($\rho=1/(44\unit{mm})=0{,}0227\unit{mm^{-1}}$).}
\end{frame}

% ============================================================================
\section{Implementación y simulaciones}
% ============================================================================

% ---- Implementación (arquitectura + pseudocódigo) ---------------------------
\begin{frame}{Implementación del motor}
\begin{columns}[T]
\begin{column}{0.52\textwidth}
\textbf{Arquitectura:}
\begin{itemize}\small
  \item Motor en Java; cada paso se computa sobre una \textbf{instantánea inmutable} de huecos y
  velocidades $\Rightarrow$ el resultado \emph{no} depende del orden de iteración.
  \item \textbf{No-solapamiento} garantizado por construcción ($d_i\le g_i+d_{i+1}$).
  \item Inserción incremental \textbf{por empuje} (llega a $N=30$ exacto).
  \item Determinista dada la realización (un único PRNG; con $p=0$ no se consume aleatoriedad).
\end{itemize}
\end{column}
\begin{column}{0.46\textwidth}
\textbf{Paso (pseudocódigo):}
{\footnotesize
\begin{block}{}
\begin{tabbing}
\quad\=\kill
por cada paso (instantánea):\\
\>R1: $\hat v_i\leftarrow\min(v_i+1,\,v_{\max,i})$\\
\>R3: con prob $p$: $\hat v_i\leftarrow\max(0,\hat v_i-1)$\\
\>R2: $d_i\leftarrow\min(\hat v_i,\,g_i+d_{i+1})$\\
\>R4: $x_i\leftarrow(x_i+d_i)\bmod L$
\end{tabbing}
\end{block}}
\end{column}
\end{columns}
\vfill
{\footnotesize Emite \textbf{solo} las variables físicas del modelo ($x$, $v$); color y observables se derivan
\emph{después}, en el análisis.}
\end{frame}

% ---- Simulaciones y observables (con definición matemática) -----------------
\begin{frame}{Simulaciones y observables}
\begin{columns}[T]
\begin{column}{0.48\textwidth}
\textbf{Parámetros:} $N\in\{5,\dots,30\}$, $p\in\{0;0{,}1;\dots;0{,}4\}$.\\[2pt]
\textbf{Protocolos:} $N$ fijo; incremental $5\to30$ (tres órdenes).\\[2pt]
\textbf{$M=30$ realizaciones} por punto; error $=$ desvío \emph{entre realizaciones}.
\end{column}
\begin{column}{0.5\textwidth}
\textbf{Observables (post-simulación):}
{\small
\[ \bar v=\frac{1}{N\,T}\sum_{t}\sum_{i=1}^{N}v_i(t) \]
\[ \sigma_{\bar v}=\sqrt{\tfrac{1}{M-1}\textstyle\sum_{m=1}^{M}(\bar v^{(m)}-\langle\bar v\rangle)^2} \]}
\begin{itemize}\small
  \item $T$: ventana \textbf{estacionaria} (corte por inspección, no \% fijo).
  \item $\rho_i=1/(\Delta x\,(\ell+g_i))$; distribuciones y diagrama fundamental.
\end{itemize}
\end{column}
\end{columns}
\end{frame}

% ============================================================================
\section{Resultados}
% ============================================================================

% ---- ANIMACIÓN (estudio N fijo) --------------------------------------------
\begin{frame}{Dinámica ($N$ fijo): flujo libre vs.\ saturación}
\begin{columns}[c]
\begin{column}{0.5\textwidth}\centering
\videohero{N5_p01_CONTACTO_PURO_FIXED_N_r1}{\linewidth}\\
{\footnotesize $N=5$ --- flujo libre}
\end{column}
\begin{column}{0.5\textwidth}\centering
\videohero{N30_p01_CONTACTO_PURO_FIXED_N_r1}{\linewidth}\\
{\footnotesize $N=30$ --- saturación (anillo a contacto)}
\end{column}
\end{columns}
\vfill\centering \playnota\ Color $=$ velocidad.
\end{frame}

% ---- observable temporal (N fijo) ------------------------------------------
\begin{frame}{Estacionario por inspección ($N$ fijo)}
\centering
\includegraphics[height=0.74\textheight]{evolucion_temporal_contacto_puro_sin_orden_fixed.png}\\
{\footnotesize El corte del transitorio se elige mirando la serie (no un \% fijo).}
\end{frame}

% ---- curva respuesta-estímulo (N fijo) -------------------------------------
\begin{frame}{Velocidad media vs.\ $N$ ($N$ fijo, por $p$)}
\centering
\includegraphics[height=0.72\textheight]{velocidad_media_vs_N_contacto_puro_sin_orden_fixed.png}\\
{\footnotesize Meseta cerca del \textbf{más lento} ($\approx90$). En $N=30$ (singular) el descenso depende de $p$;
los puntos $p\ge0{,}2$ \textbf{no son estacionarios} en el horizonte simulado (respuesta de horizonte finito). Comparable: $p=0{,}1$.}
\end{frame}

% ---- ANIMACIÓN (estudio orden) ---------------------------------------------
\begin{frame}{Efecto del orden de inserción (protocolo incremental)}
\begin{columns}[c]
\begin{column}{0.33\textwidth}\centering
\videohero{INC_p01_CONTACTO_PURO_ASCENDING_r1}{\linewidth}\\
{\footnotesize creciente}
\end{column}
\begin{column}{0.33\textwidth}\centering
\videohero{INC_p01_CONTACTO_PURO_DESCENDING_r1}{\linewidth}\\
{\footnotesize decreciente}
\end{column}
\begin{column}{0.33\textwidth}\centering
\videohero{INC_p01_CONTACTO_PURO_RANDOM_r1}{\linewidth}\\
{\footnotesize aleatorio}
\end{column}
\end{columns}
\vfill\centering \playnota\ Incremental $5\to30$, $p=0{,}1$.
\end{frame}

% ---- observable temporal (incremental) : escalera de inserción -------------
\begin{frame}{Evolución temporal incremental (escalera de inserción)}
\centering
\includegraphics[height=0.72\textheight]{evolucion_temporal_incremental_contacto_puro_p01_incremental.png}\\
{\footnotesize Cada $180\unit{s}$ entran $5$ vehículos: la velocidad media salta y se reacomoda. El observable de cada
fase se promedia sobre su ventana completa de $180\unit{s}$ (igual que el artículo).}
\end{frame}

% ---- curva respuesta-estímulo (orden) --------------------------------------
\begin{frame}{Velocidad media vs.\ $N$ por orden (\(\approx\)Fig.~2)}
\centering
\includegraphics[height=0.74\textheight]{velocidad_media_incremental_contacto_puro_p01_incremental.png}\\
{\footnotesize Decreciente $>$ aleatorio $>$ creciente a densidad baja/media; convergen a saturación.}
\end{frame}

% ---- distribuciones --------------------------------------------------------
\begin{frame}{Distribuciones por orden (\(\approx\)Figs.~3 y 4)}
\begin{columns}[c]
\begin{column}{0.5\textwidth}\centering
\includegraphics[width=\linewidth]{pdf_densidad_incremental_contacto_puro_ascending_p01_incremental.png}
\end{column}
\begin{column}{0.5\textwidth}\centering
\includegraphics[width=\linewidth]{pdf_velocidad_incremental_contacto_puro_ascending_p01_incremental.png}
\end{column}
\end{columns}
\vfill
{\footnotesize Orden creciente: la \textbf{densidad} pica en $1/(44\unit{mm})$ y la \textbf{velocidad} queda casi
independiente de $N$ (el más lento está desde $N=5$).}
\end{frame}

% ---- diagrama fundamental --------------------------------------------------
\begin{frame}{Diagrama fundamental por orden (\(\approx\)Fig.~5D)}
\centering
\includegraphics[height=0.72\textheight]{diagrama_fundamental_contacto_puro_incremental_180s_p0.1.png}\\
{\footnotesize A densidad baja/media el aleatorio queda acotado entre creciente y decreciente; a saturación convergen.}
\end{frame}

% ---- validación ------------------------------------------------------------
\begin{frame}{Validación del motor (variante clásica B)}
\begin{columns}[c]
\begin{column}{0.5\textwidth}
\begin{itemize}\small
  \item Variante B homogénea, $\ell=1$, $p=0$, reparto uniforme.
  \item Reproduce el diagrama \textbf{triangular} analítico
  \[ Q(\rho)=\min(\rho\,v_{\max},\,1-\rho) \]
  a precisión de máquina ($10^{-9}$ en tests).
  \item Contacto puro \textbf{no} se valida contra esta curva.
\end{itemize}
\end{column}
\begin{column}{0.48\textwidth}\centering
\includegraphics[width=\linewidth]{validacion_triangular.png}
\end{column}
\end{columns}
\end{frame}

% ---- SENSIBILIDAD (dt, L, Δx) ----------------------------------------------
\begin{frame}{Sensibilidad a la discretización ($\Delta x$, $\Delta t$, $L$)}
\centering
\includegraphics[width=\linewidth]{sensibilidad_dt_L_dx.png}\\[2pt]
{\footnotesize $\bar v(N)$ (contacto puro, $p=0{,}1$) variando cada constante. \textbf{$L$} irrelevante ($\le0{,}03\unit{mm/s}$);
\textbf{$\Delta x$ y $\Delta t$} sólo importan por su cociente $\Delta v=\Delta x/\Delta t$ (convergido en el espacio);
la \emph{forma} es robusta, el cuanto $\Delta v$ fija un piso de cuantización ($\le1{,}2\unit{mm/s}$ al refinar).}
\end{frame}

% ---- limitaciones ----------------------------------------------------------
\begin{frame}{Limitaciones}
\begin{itemize}
  \item No reproduce la \textbf{cola de densidad} $>1/(44\unit{mm})$ del experimento (impone no-solapamiento).
  \item $N=30$ es \textbf{singular} y depende de $L=1320$ (el experimento usa $\approx1313$).
  \item La caída bajo el más lento aparece ya a $p=0{,}1$ (satura $\approx85$, estacionario); a $p\ge0{,}2$
  el anillo lleno \textbf{no estabiliza} en el horizonte simulado (relaja lentísimo / queda casi detenido),
  régimen ajeno al experimento. Punto comparable: $p=0{,}1$.
\end{itemize}
\end{frame}

% ---- conclusiones ----------------------------------------------------------
\begin{frame}{Conclusiones}
\begin{itemize}
  \item El autómata con \textbf{contacto puro} reproduce las \emph{tendencias} del diagrama fundamental
  de VDV: velocidad casi plana y caída cerca del contacto.
  \item El \textbf{orden de inserción} modula el diagrama (aleatorio acotado a densidad baja/media; convergen a saturación).
  \item Reproduce el pico de densidad en $1/(44\unit{mm})$ y el angostamiento de la velocidad con $N$.
  \item La variante clásica \textbf{valida el motor} contra la curva triangular analítica.
  \item Los resultados son \textbf{robustos} a la discretización: $L$ es irrelevante, y $\Delta x$/$\Delta t$
  sólo importan por $\Delta v$ (calibración casi convergida).
\end{itemize}
\end{frame}

\end{document}
